\begin{thebibliography}{99}
\addcontentsline{toc}{section}{References}
\bibitem{WikiRoot}
Wikipedia contributors. \emph{Computational complexity theory}, and the selected immediate outgoing subject links recorded in \texttt{sources/coverage\_map.json}. Retrieved 22 November 2026. Used to delimit the requested subject matter, not as the authority for the arithmetic proofs.
\url{https://en.wikipedia.org/wiki/Computational_complexity_theory}.

\bibitem{ProveItPR}
Vladimir Reshetnikov, \emph{ProveIt}, revision \texttt{3eb195cab4cc3d552a9eca78880c6d428f3e44f9}. \path{Logic/PrimitiveRecursiveArithmetic/ConciseIPRA/Lean/ConciseIPRA/PR.lean}. Inspected source for \texttt{PR}, \texttt{PTm}, \texttt{PRAAx}, and quantifier-free induction.
\url{https://github.com/VladimirReshetnikov/ProveIt/blob/3eb195cab4cc3d552a9eca78880c6d428f3e44f9/Logic/PrimitiveRecursiveArithmetic/ConciseIPRA/Lean/ConciseIPRA/PR.lean}.

\bibitem{ProveItBounded}
\emph{ProveIt}, same revision. \path{Derivations/Lean/IPRADerivations/Foundations/BoundedFormulas.lean} and \path{BoundedSyntaxInduction.lean}, under \path{Logic/PrimitiveRecursiveArithmetic}. Inspected bounded-formula and object-induction interfaces.
\url{https://github.com/VladimirReshetnikov/ProveIt/blob/3eb195cab4cc3d552a9eca78880c6d428f3e44f9/Logic/PrimitiveRecursiveArithmetic/Derivations/Lean/IPRADerivations/Foundations/BoundedSyntaxInduction.lean}.

\bibitem{ProveItDiagonal}
\emph{ProveIt}, same revision. \texttt{Foundations/DiagonalClock.lean} in the \texttt{IPRADerivations} library. Object derivations \texttt{clock\_exp} and \texttt{cost\_exp}.
\url{https://github.com/VladimirReshetnikov/ProveIt/blob/3eb195cab4cc3d552a9eca78880c6d428f3e44f9/Logic/PrimitiveRecursiveArithmetic/Derivations/Lean/IPRADerivations/Foundations/DiagonalClock.lean}.

\bibitem{ProveItBaseline}
\emph{ProveIt}, same revision. \emph{Computational complexity in intuitionistic primitive recursive arithmetic}, research packet dated 20 November 2026. README, claims ledger, foundational and diagonal sections, and object-iPRA roadmap inspected.
\url{https://github.com/VladimirReshetnikov/ProveIt/tree/3eb195cab4cc3d552a9eca78880c6d428f3e44f9/Logic/PrimitiveRecursiveArithmetic/docs/proofs/foundations/ipra_complexity_foundations_2026-11-20}.

\bibitem{Cook1971}
Stephen A. Cook. The complexity of theorem-proving procedures. \emph{Proceedings of the Third Annual ACM Symposium on Theory of Computing}, pp.~151--158, 1971. DOI: 10.1145/800157.805047. Retrieved transcription of the primary paper:
\url{https://www.cs.cmu.edu/~15455/resources/Cook1971-complx-thm-proof.pdf}.

\bibitem{Karp1972}
Richard M. Karp. Reducibility among combinatorial problems. In R. E. Miller and J. W. Thatcher, eds., \emph{Complexity of Computer Computations}, pp.~85--103, 1972. DOI: 10.1007/978-1-4684-2001-2\_9. Retrieved reprint:
\url{https://www.cs.umd.edu/~gasarch/BLOGPAPERS/Karp.pdf}.

\bibitem{AspvallPlassTarjan1979}
Bengt Aspvall, Michael F. Plass, and Robert Endre Tarjan. A linear-time algorithm for testing the truth of certain quantified Boolean formulas. \emph{Information Processing Letters} 8(3), pp.~121--123, 1979. DOI: 10.1016/0020-0190(79)90002-4.
\url{https://mathweb.ucsd.edu/~sbuss/CourseWeb/Math268_2007WS/2SAT.pdf}.
Our simpler bounded algorithm proves polynomial time, not the optimal linear bound.

\bibitem{Savitch1970}
Walter J. Savitch. Relationships between nondeterministic and deterministic tape complexities. \emph{Journal of Computer and System Sciences} 4(2), pp.~177--192, 1970. DOI: 10.1016/S0022-0000(70)80006-X.
\url{https://www.sciencedirect.com/science/article/pii/S002200007080006X}.

\bibitem{Immerman1988}
Neil Immerman. Nondeterministic space is closed under complementation. \emph{SIAM Journal on Computing} 17(5), pp.~935--938, 1988. DOI: 10.1137/0217058.
\url{https://epubs.siam.org/doi/abs/10.1137/0217058}.

\bibitem{Szelepcsenyi1988}
R\'obert Szelepcs\'enyi. The method of forced enumeration for nondeterministic automata. \emph{Acta Informatica} 26, pp.~279--284, 1988. DOI: 10.1007/BF00299636.
\url{https://doi.org/10.1007/BF00299636}.

\bibitem{HartmanisStearns1965}
Juris Hartmanis and Richard E. Stearns. On the computational complexity of algorithms. \emph{Transactions of the American Mathematical Society} 117, pp.~285--306, 1965. DOI: 10.1090/S0002-9947-1965-0170805-7.
\url{https://www3.risc.jku.at/people/schreine/courses/compcomp/HartmanisStearns1965.pdf}.

\bibitem{HennieStearns1966}
F. C. Hennie and R. E. Stearns. Two-tape simulation of multitape Turing machines. \emph{Journal of the ACM} 13(4), pp.~533--546, 1966. DOI: 10.1145/321356.321362.
\url{https://www.khoury.northeastern.edu/~viola/classes/papers/HennieStearns66.pdf}.

\bibitem{ChandraKozenStockmeyer1981}
Ashok K. Chandra, Dexter C. Kozen, and Larry J. Stockmeyer. Alternation. \emph{Journal of the ACM} 28(1), pp.~114--133, 1981. DOI: 10.1145/322234.322243. Author's publication record:
\url{https://www.cs.cornell.edu/kozen/Papers/papers_by_year.htm}.

\bibitem{adleman1978}
Leonard M. Adleman. Two theorems on random polynomial time. \emph{19th Annual Symposium on Foundations of Computer Science}, pp.~75--83, 1978. DOI: 10.1109/SFCS.1978.37. Bibliographic attribution; full-paper retrieval failed. The finite proof is developed in this report and compared with the official notes below.
\url{https://doi.org/10.1109/SFCS.1978.37}.

\bibitem{lautemann1983}
Clemens Lautemann. BPP and the polynomial hierarchy. \emph{Information Processing Letters} 17(4), pp.~215--217, 1983. DOI: 10.1016/0020-0190(83)90044-3. Publisher abstract retrieved.
\url{https://www.sciencedirect.com/science/article/pii/0020019083900443}.

\bibitem{shannon1949}
Claude E. Shannon. The synthesis of two-terminal switching circuits. \emph{Bell System Technical Journal} 28(1), pp.~59--98, 1949. DOI: 10.1002/j.1538-7305.1949.tb03624.x. Official publication record:
\url{https://www.nokia.com/bell-labs/publications-and-media/publications/the-synthesis-of-two-terminal-switching-circuits/}.
The report specifies its own finite circuit basis and explicit constants.

\bibitem{katz2011}
Jonathan Katz. \emph{Notes on Complexity Theory}. University of Maryland, 2011. Official instructor notes, especially the circuit, counting, and randomized-computation lectures.
\url{https://www.cs.umd.edu/~jkatz/complexity/f11/all.pdf}.

\bibitem{AroraBarak}
Sanjeev Arora and Boaz Barak. \emph{Computational Complexity: A Modern Approach}. Cambridge University Press, 2009. Author-hosted draft, used for orientation; access limitations are recorded in the source manifest.
\url{https://www.cs.princeton.edu/~arora/pubs/webimage.pdf}.

\bibitem{mit2010sorting}
Constantinos Daskalakis. \emph{6.006 Introduction to Algorithms, Lecture 10: Sorting}. Massachusetts Institute of Technology, Fall 2010. Official course slides.
\url{https://courses.csail.mit.edu/6.006/fall10/lectures/lec10.pdf}.

\bibitem{blais2025communication}
Eric Blais. \emph{Models of Computation, Chapter 21: Communication Complexity}. University of Waterloo, Winter 2025. Official instructor notes.
\url{https://cs.uwaterloo.ca/~eblais/cs365/w25/communication}.

\bibitem{Shi2003}
Yaoyun Shi. Both Toffoli and controlled-NOT need little help to do universal quantum computation. \emph{Quantum Information and Computation} 3(1), pp.~84--92, 2003. DOI: 10.26421/QIC3.1-7. Primary preprint:
\url{https://arxiv.org/abs/quant-ph/0205115}.
The universality theorem is cited as context, not assumed as an iPRA theorem in this packet.

\bibitem{AdlemanDeMarraisHuang1997}
Leonard M. Adleman, Jonathan DeMarrais, and Ming-Deh A. Huang. Quantum computability. \emph{SIAM Journal on Computing} 26(5), pp.~1524--1540, 1997. DOI: 10.1137/S0097539795293639. Publisher abstract retrieved.
\url{https://epubs.siam.org/doi/10.1137/S0097539795293639}.

\bibitem{Babai2016}
L\'aszl\'o Babai. \emph{Graph Isomorphism in Quasipolynomial Time}. Primary preprint, arXiv:1512.03547, with subsequent author revisions. Cited only to correct an outdated baseline in the source article; the algorithm is not proved here.
\url{https://arxiv.org/abs/1512.03547}.

\bibitem{ClayPNP}
Clay Mathematics Institute. \emph{P vs NP}. Official problem statement and status, retrieved 22 November 2026.
\url{https://www.claymath.org/millennium/p-vs-np/}.
\end{thebibliography}
